\chapter*{ТЕРМИНЫ И ОПРЕДЕЛЕНИЯ}
\addcontentsline{toc}{chapter}{ТЕРМИНЫ И ОПРЕДЕЛЕНИЯ}
\setlength{\parindent}{0pt}
% Перечень располагается столбцом, слева без абзацного отступа в алфавитном порядке,
% справа через тире — определение, без знаков препинания в конце строки (п. 4.9.2)
% Термины пишутся без жирного шрифта
Автоматическая проверка --- оценивание учебных заданий программной системой без непосредственной ручной проверки преподавателем\\[0.5em]
Большая языковая модель (БЯМ) --- нейросетевая модель, генерирующая текст\\[0.5em]
Jupyter Notebook --- интерактивная вычислительная среда, объединяющая код, текстовые пояснения и результаты выполнения в одном документе\\[0.5em]
Кросс-проверка --- взаимная проверка работ обучающимися по заданным критериям\\[0.5em]
Открытый теоретический вопрос --- задание, требующее развёрнутого текстового ответа с объяснением, интерпретацией или выводом\\[0.5em]
Юнит-тест --- автоматическая проверка отдельного фрагмента программы по заранее заданным условиям\\[0.5em]